%-----------------------------------------------------------------------------------------------------------------------------------------------%
%	The MIT License (MIT)
%
%	Copyright (c) 2015 Jan Küster
%
%	Permission is hereby granted, free of charge, to any person obtaining a copy
%	of this software and associated documentation files (the "Software"), to deal
%	in the Software without restriction, including without limitation the rights
%	to use, copy, modify, merge, publish, distribute, sublicense, and/or sell
%	copies of the Software, and to permit persons to whom the Software is
%	furnished to do so, subject to the following conditions:
%	
%	THE SOFTWARE IS PROVIDED "AS IS", WITHOUT WARRANTY OF ANY KIND, EXPRESS OR
%	IMPLIED, INCLUDING BUT NOT LIMITED TO THE WARRANTIES OF MERCHANTABILITY,
%	FITNESS FOR A PARTICULAR PURPOSE AND NONINFRINGEMENT. IN NO EVENT SHALL THE
%	AUTHORS OR COPYRIGHT HOLDERS BE LIABLE FOR ANY CLAIM, DAMAGES OR OTHER
%	LIABILITY, WHETHER IN AN ACTION OF CONTRACT, TORT OR OTHERWISE, ARISING FROM,
%	OUT OF OR IN CONNECTION WITH THE SOFTWARE OR THE USE OR OTHER DEALINGS IN
%	THE SOFTWARE.
%	
%
%-----------------------------------------------------------------------------------------------------------------------------------------------%


%============================================================================%
%
%	DOCUMENT DEFINITION
%
%============================================================================%

%we use article class because we want to fully customize the page and dont use a cv template
\documentclass[10pt,A4]{article}	


%----------------------------------------------------------------------------------------
%	ENCODING
%----------------------------------------------------------------------------------------

%we use utf8 since we want to build from any machine
\usepackage[utf8]{inputenc}		

%----------------------------------------------------------------------------------------
%	LOGIC
%----------------------------------------------------------------------------------------

% provides \isempty test
\usepackage{xifthen}

%----------------------------------------------------------------------------------------
%	FONT
%----------------------------------------------------------------------------------------

% some tex-live fonts - choose your own

%\usepackage[defaultsans]{droidsans}
%\usepackage[default]{comfortaa}
%\usepackage{cmbright}
\usepackage[default]{raleway}
%\usepackage{fetamont}
%\usepackage[default]{gillius}
%\usepackage[light,math]{iwona}
%\usepackage[thin]{roboto} 

% set font default
\renewcommand*\familydefault{\sfdefault} 	
\usepackage[T1]{fontenc}

% more font size definitions
\usepackage{moresize}		


%----------------------------------------------------------------------------------------
%	PAGE LAYOUT  DEFINITIONS
%----------------------------------------------------------------------------------------

%debug page outer frames
%\usepackage{showframe}			


%define page styles using geometry
\usepackage[a4paper]{geometry}		

% for example, change the margins to 2 inches all round
\geometry{top=1.75cm, bottom=-.6cm, left=1.5cm, right=1.5cm} 	

%use customized header
\usepackage{fancyhdr}				
\pagestyle{fancy}

%less space between header and content
\setlength{\headheight}{-5pt}		


%customize entries left, center and right
\lhead{}
\chead{ \small{Ritesh Agarwal $\cdot$ Embedded Software Development Engineer $\cdot$  San Jose, CA, US $\cdot$  \textcolor{sectcol}{\textbf{riteshja88@gmail.com}}  $\cdot$ +1 312-292-7184}}
\rhead{}


%indentation is zero
\setlength{\parindent}{0mm}

%----------------------------------------------------------------------------------------
%	TABLE /ARRAY DEFINITIONS
%---------------------------------------------------------------------------------------- 

%for layouting tables
\usepackage{multicol}			
\usepackage{multirow}

%extended aligning of tabular cells
\usepackage{array}

\newcolumntype{x}[1]{%
>{\raggedleft\hspace{0pt}}p{#1}}%


%----------------------------------------------------------------------------------------
%	GRAPHICS DEFINITIONS
%---------------------------------------------------------------------------------------- 

%for header image
\usepackage{graphicx}

%for floating figures
\usepackage{wrapfig}
\usepackage{float}
%\floatstyle{boxed} 
%\restylefloat{figure}

%for drawing graphics		
\usepackage{tikz}				
\usetikzlibrary{shapes, backgrounds,mindmap, trees}


%----------------------------------------------------------------------------------------
%	Color DEFINITIONS
%---------------------------------------------------------------------------------------- 

\usepackage{color}

%accent color
\definecolor{sectcol}{RGB}{255,150,0}

%dark background color
\definecolor{bgcol}{RGB}{110,110,110}

%light background / accent color
\definecolor{softcol}{RGB}{225,225,225}


%============================================================================%
%
%
%	DEFINITIONS
%
%
%============================================================================%

%----------------------------------------------------------------------------------------
% 	HEADER
%----------------------------------------------------------------------------------------

% remove top header line
\renewcommand{\headrulewidth}{0pt} 

%remove botttom header line
\renewcommand{\footrulewidth}{0pt}	  	

%remove pagenum
\renewcommand{\thepage}{}	

%remove section num		
\renewcommand{\thesection}{}			

%----------------------------------------------------------------------------------------
% 	ARROW GRAPHICS in Tikz
%----------------------------------------------------------------------------------------

% a six pointed arrow poiting to the left
\newcommand{\tzlarrow}{(0,0) -- (0.2,0) -- (0.3,0.2) -- (0.2,0.4) -- (0,0.4) -- (0.1,0.2) -- cycle;}	

% include the left arrow into a tikz picture
% param1: fill color
%
\newcommand{\larrow}[1]
{\begin{tikzpicture}[scale=0.58]
	 \filldraw[fill=#1!100,draw=#1!100!black]  \tzlarrow
 \end{tikzpicture}
}

% a six pointed arrow poiting to the right
\newcommand{\tzrarrow}{ (0,0.2) -- (0.1,0) -- (0.3,0) -- (0.2,0.2) -- (0.3,0.4) -- (0.1,0.4) -- cycle;}

% include the right arrow into a tikz picture
% param1: fill color
%
\newcommand{\rarrow}
{\begin{tikzpicture}[scale=0.7]
	\filldraw[fill=sectcol!100,draw=sectcol!100!black] \tzrarrow
 \end{tikzpicture}
}



%----------------------------------------------------------------------------------------
%	custom sections
%----------------------------------------------------------------------------------------

% create a coloured box with arrow and title as cv section headline
% param 1: section title
%
\newcommand{\cvsection}[1]
{
\colorbox{sectcol}{\mystrut \makebox[1\linewidth][l]{
\larrow{bgcol} \hspace{-8pt} \larrow{bgcol} \hspace{-8pt} \larrow{bgcol} \textcolor{white}{\textbf{#1}}\hspace{4pt}
}}\\
}

%create a coloured arrow with title as cv meta section section
% param 1: meta section title
%
\newcommand{\metasection}[2]
{
\begin{tabular*}{1\textwidth}{p{2.4cm} p{11cm}}
\larrow{bgcol}	\normalsize{\textcolor{sectcol}{#1}}&#2\\[12pt]
\end{tabular*}
}

%----------------------------------------------------------------------------------------
%	 CV EVENT
%----------------------------------------------------------------------------------------

% creates a stretched box as cv entry headline followed by two paragraphs about 
% the work you did
% param 1:	event time i.e. 2014 or 2011-2014 etc.
% param 2:	event name (what did you do?)
% param 3:	institution (where did you work / study)
% param 4:	what was your position
% param 5:	some words about your contributions
%
\newcommand{\cvevent}[5]
{
\vspace{8pt}
	\begin{tabular*}{1\textwidth}{p{2.3cm}  p{10.8cm} x{3.9cm}}
 \textcolor{bgcol}{#1}& \textbf{#2} & \vspace{2.5pt}\textcolor{sectcol}{#3}

	\end{tabular*}
\vspace{-12pt}
\textcolor{softcol}{\hrule}
\vspace{6pt}
	\begin{tabular*}{1\textwidth}{p{2.3cm} p{14.4cm}}
&		 \larrow{bgcol}  #4\\[3pt]
&		 \larrow{bgcol}  #5\\[6pt]
	\end{tabular*}

}

\newcommand{\cveventstart}[3]
{
\vspace{8pt}
	\begin{tabular*}{1\textwidth}{p{2.3cm}  p{10.8cm} x{3.9cm}}
 \textcolor{bgcol}{#1}& \textbf{#2} & \vspace{2.5pt}\textcolor{sectcol}{#3}

	\end{tabular*}
\vspace{-12pt}
\textcolor{softcol}{\hrule}
\vspace{6pt}
	\begin{tabular*}{1\textwidth}{p{2.3cm} p{14.4cm}}
}
\newcommand{\cveventi}[1]
{
&		 \larrow{bgcol}  #1\\[3pt]
}
\newcommand{\cveventend}
{
	\end{tabular*}
}

% creates a stretched box as 
\newcommand{\cveventmeta}[2]
{
	\mbox{\mystrut \hspace{87pt}\textit{#1}}\\
	#2
}

%----------------------------------------------------------------------------------------
% CUSTOM STRUT FOR EMPTY BOXES
%----------------------------------------- -----------------------------------------------
\newcommand{\mystrut}{\rule[-.3\baselineskip]{0pt}{\baselineskip}}

%----------------------------------------------------------------------------------------
% CUSTOM LOREM IPSUM
%----------------------------------------------------------------------------------------
\newcommand{\lorem}
{Lorem ipsum dolor sit amet, consectetur adipiscing elit. Donec a diam lectus.}



%============================================================================%
%
%
%
%	DOCUMENT CONTENT
%
%
%
%============================================================================%
\begin{document}


%use our custom fancy header definitions
\pagestyle{fancy}	


%---------------------------------------------------------------------------------------
%	TITLE HEADLINE
%----------------------------------------------------------------------------------------
\vspace{-20.55pt}

% use this for multiple words like working titles etc.
%\hspace{-0.25\linewidth}\colorbox{bgcol}{\makebox[1.5\linewidth][c]{\hspace{46pt}\HUGE{\textcolor{white}{\textsc{Ritesh Agarwal}} } \textcolor{sectcol}{\rule[-1mm]{1mm}{0.9cm}} \parbox[b]{5cm}{   \large{ \textcolor{white}{{IT Consultant}}}\\
% \large{ \textcolor{white}{{Resume}}}}
%}}

% use this for single words, e.g. CV or RESUME etc.
\hspace{-0.25\linewidth}\colorbox{bgcol}{\makebox[1.5\linewidth][c]{\HUGE{\textcolor{white}{\textsc{Ritesh Agarwal}} } \textcolor{sectcol}{\rule[-1mm]{1mm}{0.9cm}} \HUGE{\textcolor{white}{\textsc{Resume}} } }}


%----------------------------------------------------------------------------------------
%	HEADER IMAGE
%----------------------------------------------------------------------------------------

\begin{figure}[H]
\begin{flushright}
	\includegraphics[width=0.2\linewidth]{qrcode-ritesh-resume-main.pdf-github-pdf.png}
\end{flushright}
\end{figure}

%---------------------------------------------------------------------------------------
%	QR CODE (optional)
%----------------------------------------------------------------------------------------
%\vspace{-136pt}
%\hspace{0.75\linewidth}
%\includegraphics[width=103pt]{qrcode}
%\normalsize
%\vspace{88pt}

%---------------------------------------------------------------------------------------
%	META SECTION
%----------------------------------------------------------------------------------------

\vspace{-114pt}

\metasection{Status:}{Cisco Systems, Embedded Software Development Engineer}
\metasection{Fields:}{Software Development, Cloud Computing, Networking, Storage, Linux Kernel, Programming}
\metasection{Prefers:}{C, C++, Python, Linux/UNIX based host and target OS, Git / GitHub, OpenSource}
\metasection{Activities:}{Hackathons, Technical Workshops, Geek Meetups}

\vspace{6pt}

%---------------------------------------------------------------------------------------
%	SUMMARAY (optional)
%----------------------------------------------------------------------------------------

%\cvsection{Summary}\\
%Digital media graduate with four years project experience in the field of technology based assessment. Specialized in development of test-scenario engines and innovative, rich media item formats. Master studies focused on teams from different disciplines and cultural backgrounds on solutions for complex problems.  Prior knowledge has been collected in he field of usability / accessibility during bachelor studies.\\

%============================================================================%
%
%	CV SECTIONS AND EVENTS (MAIN CONTENT)
%
%============================================================================%

%---------------------------------------------------------------------------------------
%	EXPERIENCE
%----------------------------------------------------------------------------------------
\cvsection{EXPERIENCE}

%March 2015 - Present
\cveventstart{2015 / 03 }{Software Development Engineer}{Cisco Systems}
\cveventi{\textbf{Technology Area:} NFV, Virtualization, Virtual Packet Core Platform, DPDK}
\cveventi{Debugging support for DPDK based critical process.}
\cveventi{Involved in deployment of various combinations of the NFV deployments for both KVM and Vmware ESXi.}
\cveventi{Development of features and Bug fixing support.}
\cveventi{Debugging features : minicore of a heavy weight multithreaded process.}
\cveventi{Bulkstats addition/modification and integrate it with the Bulkstats ecosystem.}
\cveventi{Work in areas of SRIOV, Passhthrough, Virtio NICs in KVM.}
\cveventi{32 bit to 64 bit conversion of forwarder task.}
\cveventend
%\textcolor{softcol}{\hrule}

%(May 2014 – May 2015).
\cveventstart{2014 - 2015}{Research Assistant}{KernelSec Lab, UIC}
\cveventi{\textbf{Technology Area:} Xen Hypervisor, Operating Systems, Asynchronous calls, Events.}
\cveventi{\textbf{Project:} IPC Subsystem Cleanups}
\cveventi{Analyse how to cleanup incomplete events that gets posted in the kernel by system calls in the IPC subsystem area.}
\cveventi{Identify test scenarios and write test cases for those scenarios.}
\cveventi{Provide a way to handle incomplete events in the event of Program termination normally or forcefully.}
\cveventi{Publish report in latex for the work done.}
\cveventend
%\textcolor{softcol}{\hrule}

%
%(June 2013 – August 2013)
\cveventstart{2013}{Graduate Summer Intern}{Yahoo! Inc}
\cveventi{\textbf{Technology Area:} OpenStack, Cloud Computing, Nicira, Software Defined Networking(SDN), Open vSwitch , REST APIs, Web Services, Python, Python Request Library.}
\cveventi{\textbf{Project: Monitoring Solution for Nicira NVP / Openstack Networking.}}
\cveventi{Collect network related statistics from Nicira NVP via Web APIs and Open vSwitch via command line and parse them to be able to facilitate easy analysis of networking related problems for Virtual Machines on the cloud. Integrate with the Internal Monitoring Cloud called YAMAS that mantained time-series and helped plot graphs.}
\cveventi{The project mainly involved understanding the virtual network topology and getting well versed with SDN Concepts and then come up with a feasible solution as to of what information can be collected from Nicira and OVS and will be useful to facilitate  monitoring.}
\cveventend
%\textcolor{softcol}{\hrule}

\newpage
%Jan 2013 - May 2014 (except Summer 2013)
\cveventstart{2013 - 2014}{Graduate Research Assistant}{RRC, UIC}
\cveventi{Job Responsibilities involved variety of work in different areas.}
\cveventi{Administration for Blade Servers, ZFS Servers, Samba Server, Git server, Mail server, Web Servers, Database Server and Network.}
\cveventi{Provide support and suggest solutions to researchers/professors for different needs.}
\cveventend
%\textcolor{softcol}{\hrule}

%July 2010 - August 2012 
\cveventstart{2010 - 2012}{Software Development Engineer}{Cisco Systems}
\cveventi{\textbf{Technology Area: 3G / UMTS, GTP (GPRS Tunneling Protocol), 4G / LTE Packet core networks, CISCO ASR5000, Carrier Grade Operating System(StarOS), Network Applications.
}}
%Job Responsibilities
\cveventi{Implement new features in product, CLI Approval, CLI Announcement, Requirement Document.}
\cveventi{Validate \ fix issues reported by customer / QA team.}
\cveventi{Browse through huge log files to explore the correct code path.}
\cveventi{Load, analyse core dumps in gdb, analyse linked lists in core dumps.}
\cveventi{Unit Test cases in Python Framework.}
\cveventi{Write bash, expect scripts to come up with complex network setups.}
\cveventi{\textbf{Achievements at Work :}Developed VIM-CDETS a vim interface for CDETS(Bug Tracking Tool).}
\cveventend

%\textcolor{softcol}{\hrule}





%	EDUCATION SECTION
%--------------------------------------------------------------------------------------
\cvsection{EDUCATION}
%Fall 2012 to Summer 2015
\cveventstart{May 2015}{MS in Computer Science}{Univ. of Illinois, Chicago}
\cveventi{Cumulative GPA 3.85 / 4.0}
\cveventi{\textbf{Courses in Areas:} Computer Networks, Computer Architecture, System Security, Object Oriented Development, Algorithms, Parallel Processing, Applied Graph Theory.}
\cveventi{\textbf{Projects in Networking area:} HTTP Client/Server(threaded), DNS Resolver, Bittorrent(event based), Reliable Transport over a non-reliable channel(TCP), Study of Internet via distributed network hosts.}
\cveventi{\textbf{Projects in Microprocessors and Computer Architecture area:} Cache replacement policy for Simplescalar (Access pattern analysis), Study of Compiler optimization technique for matrix multiplication using blocking method to improve cache use.}
\cveventi{\textbf{Projects in Object Oriented Languages and Environment:} Object Oriented Design with C++, Memory leak analysis with valgrind, Cincom Smalltalk, iOS App Development with Xcode in Objective C.}
\cveventi{\textbf{Projects in Parallel Processing:} MPI Programming.}
\cveventend

%\textcolor{softcol}{\hrule}

%
\cveventstart{May 2010}{BE in Information Technology}{P.I.C.T, Univ. of Pune}
\cveventi{Cumulative GPA 3.60 / 4.0}
\cveventend

%\textcolor{softcol}{\hrule}

%
\cveventstart{2010}{Project: Reconfigurable Virtual Storage Device}{P.I.C.T, Univ. of Pune}
\cveventi{Guided by Ratnadeep Joshi (Toshiba) and Furquan Shaikh (Google).}
\cveventi{Undergraduate Research Project in Linux kernel.}
\cveventi{Ranked 3rd in Impetus and Concepts 2010 in Database and Storage Systems Area.}
\cveventi{Stack Device driver that combined SSD and Hard Drive and exported itself as a virtual block storage device.}
\cveventi{Processed bio requests from filesystems and mapped it to a bio on the respective device.}
\cveventi{Heuristic Access pattern analysis to characterize current hot blocks using EWMA.}
\cveventi{Workqueues for block relocation.}
\cveventi{Journaling of block relocation sub-operations for recovery from incomplete operations due to system failure.}
\cveventi{Performance: Achieved the performance close to that of SSD, with a SSD contributing to 30\% of to the total storage.}
\cveventend

\newpage
%\textcolor{softcol}{\hrule}

%
\cveventstart{2007}{Antivirus}{P.I.C.T, Univ. of Pune}
\cveventi{Won First Prize in State Level Project Competition in 2007.}
\cveventi{A miniproject developed in C and a GUI with DOS Mode Graphics, C++ for GUI Library.}
\cveventi{Mouse Handling using Interrupts via Assembly Language code.}
\cveventi{Scan for virus signatures in files and detect infected files.}
\cveventi{Reverse the effects of virus that were non-destructive in nature.}
\cveventend

%\textcolor{softcol}{\hrule}

%
\cveventstart{2008}{PortScanner}{P.I.C.T, Univ. of Pune}
\cveventi{Distributed Port Scanning by exploiting the ID field in the IP Header.}
\cveventi{Worked much faster and was almost correct with its results.}
\cveventend

%---------------------------------------------------------------------------------------
%	TECHNICAL SKILLS
%----------------------------------------------------------------------------------------
\cvsection{TECHNICAL SKILLS}
\cveventstart{}{PROGRAMMING LANGUAGES KNOWN:}{}
\cveventi{C(Proficient), C++, Assembly Language(8086), Python}
\cveventend

\cveventstart{}{COMPUTER SKILLS / TOOLS KNOWN:}{}
\cveventi{\textbf{Code Browsing / Refactoring :} vim (also developed a plugin), cscope, ctags, awk, sed.}
\cveventi{\textbf{Debugging Program / Cores:}}
\cveventi{\textbf{Gdb(GNU Debugger):} used conditional breakpoints and watchpoints to narrow down the problem, analysed core file, stack frames, inspected code paths, examined packets(hexdump) in memory, wrote gdb scripts / helpers for e.g to traverse linked list.}
\cveventi{\textbf{Pdb(Python Debugger):} Debugged python scripts.}
\cveventi{\textbf{Network Tools :} wireshark and tcpdump(analyzed packet dumps), hping3 and ip (tailored packets), callgen and SAE-SIM(to tailor GTP packets).}
\cveventi{\textbf{Programming :} Socket programming, Device drivers, kernel modules, system calls.}
\cveventi{\textbf{Log Analysis :} NPU Manager logs / packet-dumps, Session Manager logs.}
\cveventi{\textbf{Web :} curl, Python Request Library to use REST APIs}
\cveventi{\textbf{Scripting :} bash, expect, GNU Screen.}
\cveventi{\textbf{Versioning System :} BitKeeper, svn, git.}
\cveventi{\textbf{Bug Tracking Tools :} CDETS, GNATS, Bugzilla.}
\cveventi{\textbf{Kernel Hacking:} Academic level idea of with hands-on in ProcFS, DebugFS, Tracing, Kprobes, Workqueues, BlockIO, the / mutex.}
\cveventi{\textbf{Optimization :} BRANCH\_PREDICT, Likely / Unlikely.}
\cveventi{\textbf{Publishing :} latex, gnuplot.}
\cveventi{\textbf{Virtualization :} KVM / Qemu, VMware ESXi, virtio, SR-IOV, PCI-PassThrough.}
\cveventend

\cveventstart{}{OTHER KEYWORDS:}{}
\cveventi{OpenVSwitch, Cloud Computing, Storage / Network QoS, Solid State Drives, Operating Systems, Filesystems, UNIX, Microprocessors, Regular Expressions, TCP / IP, IPv6, P2P(churn, co-operative tit-for-tat), Telecommunications, Embedded System.}
\cveventend

\cveventstart{}{EXTRA-CURRICULAR}{}
\cveventi{Won numerous State and National Level Programming, Technical Quiz and Hacking Competitions.}
\cveventend




%-------------------------------------------------------------------------------------------------
%	ARTIFICIAL FOOTER (fancy footer cannot exceed linewidth) 
%--------------------------------------------------------------------------------------------------

\null
\vspace*{\fill}
\hspace{-0.25\linewidth}\colorbox{bgcol}{\makebox[1.5\linewidth][c]{\mystrut \small \textcolor{white}{www.linkedin.com/in/riteshja88} $\cdot$ \textcolor{white}{github.com/riteshja88}}}




%============================================================================%
%
%
%
%	DOCUMENT END
%
%
%
%============================================================================%
\end{document}
